% Three ways to describe a type, one schema, one SDL document. Referenced from
% chapter 02. Bare tikzpicture: the figure environment, caption and label live
% at the call site. Uses only arrows.meta and positioning (loaded in packages.tex).
\begin{tikzpicture}[
    font=\small,
    src/.style={draw, rounded corners=2pt, minimum width=40mm,
                minimum height=17mm, align=center, font=\scriptsize},
    odd/.style={draw, dashed, rounded corners=2pt, minimum width=40mm,
                minimum height=17mm, align=center, font=\scriptsize},
    core/.style={draw, thick, rounded corners=2pt, minimum width=34mm,
                 minimum height=12mm, align=center},
    emit/.style={draw, rounded corners=2pt, minimum width=32mm,
                 minimum height=12mm, align=center},
    flow/.style={-{Stealth[length=2mm]}, thick},
    note/.style={font=\scriptsize\itshape, align=center}
  ]

  \node[src] (impl) at (-4.8,2.5)
    {{\small Implementation-first}\\[2pt]
     C\# types plus attributes,\\
     read by a source generator};

  \node[src] (code) at (-4.8,0)
    {{\small Code-first}\\[2pt]
     descriptor classes:\\
     \texttt{ObjectType<T>} and\\
     \texttt{Configure}};

  \node[odd] (schemafirst) at (-4.8,-2.5)
    {{\small Schema-first}\\[2pt]
     an SDL document plus\\
     runtime type bindings};

  \node[core] (schema) at (0,0)   {one schema object\\in memory};
  \node[emit] (sdl)    at (4.5,0) {the same SDL,\\byte for byte};

  \foreach \s in {impl,code,schemafirst}{%
    \draw[flow] (\s) -- (schema);
  }
  \draw[flow] (schema) -- (sdl);

  \node[note, anchor=south] at (-4.8,3.45)
    {Mosaic uses implementation-first};
  \node[note, anchor=north] at (-4.8,-3.5)
    {dashed: supported, but absent\\from the v16 documentation};

\end{tikzpicture}
